\selectlanguage{ngerman}%
\frenchspacing%
\begin{ledgroupsized}[r]{120mm}%
\footnotesize%
\pstart%
\noindent\textbf{Überlieferung:}%
\pend%
\end{ledgroupsized}%
\begin{ledgroupsized}[r]{114mm}%
\footnotesize%
\pstart%
\parindent -6mm%
\makebox[6mm][l]{\textit{L}}%
Notiz:
LH XXXV 8,\;30 Bl\;121.
Ein Blatt 2\textsuperscript{o},
Ränder beschnitten;
Wasserzeichen: %
Pariser Papier.\indexO{Paris}
Sechs ganzseitig beschriebene Zeilen in der oberen Hälfte von Bl.\;\href{\Scan{LH_35_08_30_121r}}{121\,r\textsuperscript{o}}.
Überliefert sind dort zusätzlich die im vorliegenden Band edierte Notiz N.\;\ref{Cc2_847-2}, die auf Bl.\;\href{\Scan{LH_35_08_30_121r}}{121\,r\textsuperscript{o}} unmittelbar nach N.\;\ref{Cc2_847-1} folgt, sowie umfangreichere mathematische Aufzeichnungen, die den Rest von Bl.\;\href{\Scan{LH_35_08_30_121r}}{121\,r\textsuperscript{o}} füllen; diese werden voraussichtlich in einem künftigen Band von \textit{LSB} VII ediert.
Zwei Textabschnitte von Bl.\;\href{\Scan{LH_35_08_30_121r}}{121\,r\textsuperscript{o}} werden auf die sonst unbeschriebene Hinterseite (Bl.\;\href{\Scan{LH_35_08_30_121v}}{121\,v\textsuperscript{o}}) fortgesetzt.
\newline%
Cc~2, Nr.\;847
\pend%
\end{ledgroupsized}%
\vspace{5mm}%
\begin{ledgroup}%
\footnotesize%
\pstart%
\noindent%
\textbf{Datierungsgründe:}
Die Datierung der Notiz N.\;\ref{Cc2_847-1} beruht auf Eigenschaften des Textträgers sowie auf inhaltlichen Betrachtungen.
Siehe die ausführliche Begründung auf S.\;\pageref{Cc2_847-1-dat}.%
\pend%
\end{ledgroup}%
\selectlanguage{latin}%
\frenchspacing%
\vspace{8mm}%
\normalsize%
\pstart%
\noindent%
\lbrack%
\href{\Scan{LH_35_08_30_121r}}%
{121\,r\textsuperscript{o}}%
\rbrack\ %
Calculus Elasticus%
\indexS{calculus elasticus}
et
\edtext{%
similes.
Spatia%
}{%
\lemma{similes.}%
\Bfootnote{%
\textit{(1)}~Vires %
\textit{(2)}~Ce %
\textit{(3)}~Spatia%
~\textit{L}%
}}
percursa%
\indexS{spatium percursum}
in duplicata ratione temporum%
\indexS{ratio temporum duplicata}
vel celeritatum;%
\indexS{ratio celeritatum duplicata}
celeritates\indexS{celeritas}
in subduplicata spatiorum.%
\indexS{ratio spatiorum subduplicata}
\pend%
\pstart%
Embolus \textit{E}%
\indexS{embolus}
pellatur ex \textit{A}
\edtext{%
in \textit{B},
et aer \textit{AB}%
}{%
\lemma{in}%
\Bfootnote{%
\textit{(1)}~\textit{A} %
\textit{(2)}~\textit{B}. %
\textit{(a)}~ae %
\textit{(b)}~et aer \textit{AB}%
~\textit{L}%
}}
adigatur intra aerem \textit{BF}.
Porro pergendo \textit{A}
et \textit{BF} jam compressus
redigendus est intra \textit{CF},
videtur resistentia haec nova%
\indexS{resistentia aeris ad compressionem}
esse ad priorem
in composita ratione,%
\indexS{ratio composita}
\edtext{%
tum \textit{BF}%
}{%
\lemma{tum}%
\Bfootnote{%
\textit{(1)}~reciproca %
\textit{(2)}~\textit{AF} %
\textit{(3)}~\textit{BF}%
~\textit{L}%
}}
ad \textit{CF}
(quia spatium nunc minus)
tum ratione
\edtext{%
aeris nunc comprimendi ad priorem.%
}{%
\lemma{aeris}%
\Bfootnote{%
\textit{(1)}~nunc ad prio %
\textit{(2)}~nunc comprimendi ad priorem%
~\textit{L}%
}}
Videntur haec non satis consentire
\edlabel{LH_35_08_30_121-1_anspiel-1}%
\edtext{%
alibi dictis,%
}{%
\lemma{alibi dictis}%
\Cfootnote{%
Wohl Anspielung auf Leibnizens Entwürfe \textit{Schediasma de calculo elastico} (\textit{LSB} VIII,\;1 N.\;53)\cite{02617} bzw. \textit{Calculus elasticus} (\textit{LSB} VIII,\;1 N.\;54),\cite{02609} beide eigenhändig auf Dezember 1674 datiert.%
}}%
\edlabel{LH_35_08_30_121-1_anspiel-2}
quasi tantum servaretur ratio spatiorum reciproca,%
\indexS{ratio spatiorum reciproca}
\edtext{%
itaque quaerenda conciliatio,%
\indexS{conciliatio quaerenda}%
}{%
\lemma{itaque}%
\Bfootnote{%
\textit{(1)}~res %
\textit{(2)}~quaerenda conciliatio%
~\textit{L}%
}}
aut veritas%
\indexS{veritas}
an illustrabilis vis Elastica%
\indexS{vis elastica}
ea vi,
qua globus eccentricus%
\indexS{globus eccentricus in liquore gyrato}
in liquore gyrato positus%
\indexS{liquor gyratus}
in medium repertur.
\pend%
\begin{edfigure}
\includegraphics[scale=1]{%
gesamttex/edit_VIII-4/images/LH_35_08_30_121-1_d_121r.pdf}
\edcaption{\lbrack\textit{Fig.\:1}\rbrack}
\end{edfigure}
\newpage%
\vspace{3.0em}%
\pstart%
\noindent%
\centering%
*     *     *
\pend%
\vspace{3.0em}%
\selectlanguage{german}%
\frenchspacing%
\footnotesize%
\pstart%
\noindent%
\textls{Ausführliche Datierungsgründe:}%
\pagelabel{Cc2_847-1-dat}
Die undatierte Notiz N.\;\ref{Cc2_847-1} ist auf demselben Textträger überliefert wie die im vorliegenden Band edierte Notiz N.\;\ref{Cc2_847-2}, die auf Bl.\;\href{\Scan{LH_35_08_30_121r}}{121\,r\textsuperscript{o}} nach N.\;\ref{Cc2_847-1} folgt, und weitere mathematische Aufzeichnungen, die den Rest von Bl.\;\href{\Scan{LH_35_08_30_121r}}{121\,r\textsuperscript{o}} füllen und voraussichtlich in einem künftigen Band von \textit{LSB}~VII ediert werden.
Diese gemeinsam überlieferte Textgruppe lässt sich auf \textit{Dezember 1674 bis 1675~(?)} datieren.
Ausschlaggebende Gründe dafür sind (1)~das die Pariser Herkunft des Papiers bezeugende Wasserzeichen sowie (2)~die in der Notiz N.\;\ref{Cc2_847-1} (\refpassage{LH_35_08_30_121-1_anspiel-1}{LH_35_08_30_121-1_anspiel-2}) anzutreffende Anspielung auf die nach Leibnizens Angabe in Dezember 1674 abgefassten Entwürfe \textit{Schediasma de calculo elastico} (\textit{LSB} VIII,\;1 N.\;53)\cite{02617} bzw. \textit{Calculus elasticus} (\textit{LSB} VIII,\;1 N.\;54);\cite{02609}
siehe für die ausführlichen Datierungsgründe die Notiz N.\;\ref{Cc2_847-2} im vorliegenden Band.
Die Umstände der Überlieferung legen es nahe, die besagte Datierung in erster Linie für die Notiz N.\;\ref{Cc2_847-1} in Anspruch zu nehmen.
\pend%
\selectlanguage{latin}%
\frenchspacing%
\normalsize%
